\documentclass[11pt,a4paper]{article}
\usepackage[utf8]{inputenc}
\usepackage[T1]{fontenc}
\usepackage{lmodern}
\usepackage[french]{babel}
\usepackage[margin=2.3cm]{geometry}
\usepackage{amsmath,amssymb}
\usepackage{graphicx}
\usepackage{booktabs}
\usepackage{xcolor}
\usepackage{enumitem}
\usepackage{tcolorbox}
\usepackage{hyperref}
\hypersetup{colorlinks=true,linkcolor=blue!50!black,urlcolor=blue!50!black}

\newcommand{\dd}{\mathrm{d}}
\newcommand{\kd}{k_1^{d}}
\newcommand{\uc}{u_c}
\newcommand{\dW}{\delta^2\mathcal W}
\newtcolorbox{alamain}{colback=yellow!8,colframe=orange!70!black,title=À faire à la main,fonttitle=\bfseries}
\newtcolorbox{idee}{colback=blue!5,colframe=blue!50!black,title=L'idée à retenir,fonttitle=\bfseries}

\title{Localisation de l'endommagement et stabilité par le travail du second ordre\\
\large Le ressort endommageable relu avec Bažant}
\author{Note complémentaire, dépôt \texttt{1D\_damage\_spring\_system}}
\date{\today}

\begin{document}
\maketitle

\begin{center}
\small Cette note complète \texttt{explication.tex}. Première partie : dans le modèle à deux
ressorts, l'endommagement est localisé dans le ressort 1 par construction ; qu'est-ce que
cela signifie par rapport à la localisation telle que l'entendent Bažant, de Borst ou Acary et
Monerie ? Seconde partie : Chaboche et Feyel écrivent que l'étude de stabilité passe par le
travail du second ordre, en citant Bažant ; on détaille ce critère et ce qu'il dit de notre
cas. Les grandeurs sont celles de \texttt{parameters.py}.
\end{center}

\tableofcontents

\part{La localisation}

\section{Deux sens du mot « localisation »}

Dans une éprouvette réelle en traction, la déformation est d'abord uniforme. Au-delà du pic
de contrainte, elle cesse de l'être : une zone étroite continue de s'allonger et de
s'endommager pendant que le reste de l'éprouvette décharge élastiquement. On appelle cela la
\emph{localisation} de l'endommagement, et la zone qui localise finit par devenir la fissure.

Il faut distinguer deux usages de ce mot.
\begin{itemize}[nosep]
  \item La localisation comme \textbf{résultat} : on part d'un milieu uniforme, on applique la
        loi adoucissante partout, et on montre que l'état uniforme devient instable au pic et
        que la déformation se concentre spontanément. C'est l'analyse de Bažant (§13.2),
        de Borst (§3) et de la note de cours \emph{Loi de comportement élastique-endommageable}
        (poutre non uniforme).
  \item La localisation comme \textbf{hypothèse} : on décide à l'avance que l'endommagement
        est confiné dans une zone donnée, et on modélise le reste comme élastique. C'est la
        zone cohésive d'Acary et Monerie, l'interface de Chaboche et Feyel, et c'est notre
        modèle à deux ressorts.
\end{itemize}
Notre ressort 1 est la zone qui localise, notre ressort 2 est tout le reste qui décharge.
Bažant le dit lui-même : le modèle de couplage en série « typifies the behavior of many
structures » et il est le prototype le plus simple de la localisation. Mais nous avons
mis la localisation dans la géométrie dès le départ : nous ne l'observons pas, nous la
postulons.

\section{Ce que fait le modèle à un ressort}

\subsection{Le trajet stable de Bažant, et l'étape que nous sautons}
Reprenons le raisonnement de Bažant (§13.2, « Bifurcation and stable path »). Soit une
éprouvette de longueur $L$ chargée à travers un ressort de machine $C_s$, en état uniforme
dans la phase adoucissante. Pour un incrément de déplacement imposé $\delta q_1$, deux trajets
sont possibles à partir de chaque état post-pic :
\begin{itemize}[nosep]
  \item trajet 1, \emph{uniforme} : toute l'éprouvette continue de s'adoucir avec le module
        tangent $E_t < 0$ ; la raideur globale est
        $C^{(1)} = \bigl[C_s^{-1} + L/(E_tA)\bigr]^{-1}$ ;
  \item trajet 2, \emph{localisé} : un segment de longueur $h$ s'adoucit avec $E_t$, le reste
        $L-h$ décharge avec $E_u > 0$ ; la raideur globale est
        $C^{(2)} = \bigl[C_s^{-1} + h/(E_tA) + (L-h)/(E_uA)\bigr]^{-1}$.
\end{itemize}
Le travail du second ordre le long d'un trajet est $\dW = \tfrac12C\,\delta q_1^2$, et le trajet
qui se produit est celui qui le minimise (section~\ref{sec:trajet}), donc celui du plus petit
$C$. Comme $E_u > 0$ et $E_t < 0$, on a toujours $C^{(2)} < C^{(1)}$ : \emph{la déformation
localise dès que la localisation est possible, c'est-à-dire dès le pic}. Et parmi les trajets
localisés, le plus petit $C$ est obtenu pour le plus petit $h$ : la zone adoucissante est la
plus courte possible.

Notre modèle saute cette étape. Il postule $h = L_1$ (le ressort 1 tout entier s'adoucit) et
$L - h = 0$ (rien dans le béton ne décharge ; seul le ressort 2 décharge). Il n'y a plus de
choix de trajet : avec un seul degré de liberté, la bifurcation n'existe pas, il ne reste que
la perte de stabilité de l'état au point limite.

\subsection{L'équation (13.2.5) de Bažant et notre cas}
Bažant écrit la condition de stabilité de l'état uniforme, sous déplacement piloté, pour le
système éprouvette + machine :
\begin{equation}
  \text{stable si}\quad -\frac{E_t}{E_u} < \frac{1}{(L/h) - 1 + \dfrac{AE_u}{C_sh}} .
  \tag{13.2.5}
\end{equation}
Trois longueurs y interviennent : $L$, la longueur totale ; $h$, la zone qui s'adoucit ;
$AE_u/C_s$, la longueur de béton qui aurait la même souplesse que la machine. Lecture :
\begin{itemize}[nosep]
  \item à $h$ fixé, plus $L$ est grand, plus le membre de droite est petit, donc plus le module
        tangent critique $|E_t^{\text{cr}}|$ est petit : l'instabilité survient plus tôt après le
        pic. Une éprouvette longue est moins ductile ;
  \item plus la machine est souple ($C_s$ petit), même effet ;
  \item à $L$ fixé, plus $h$ est petit, plus $|E_t^{\text{cr}}|$ est petit, et
        $\lim_{h\to0}E_t^{\text{cr}} = 0$ : si la zone adoucissante peut être arbitrairement mince,
        l'instabilité survient \emph{au pic} et l'énergie dissipée est nulle.
\end{itemize}
Notre modèle est le cas $h = L = L_1$, $C_s = k_2$ : le membre de droite devient $C_sL_1/(AE_u)$,
et la condition se réduit à $-E_tA/L_1 < k_2$, c'est-à-dire $-\dd F_1/\dd u_1 < k_2$, soit $k_t > 0$.
Nous retrouvons notre critère, mais nous avons perdu les deux premiers effets : la longueur
de béton qui décharge et la taille de la zone de rupture ne sont plus des paramètres.

\section{Ce que l'hypothèse coûte et ce qu'elle épargne}

\subsection{Ce qu'elle coûte}
\begin{itemize}[nosep]
  \item On ne voit pas la localisation se produire ; on ne peut pas demander où ni sur quelle
        longueur elle a lieu.
  \item La ductilité apparente est fixée par $L_1$. À $f_t$ et $G_f$ donnés, la pente de la
        descente de $F_1(u_1)$ est d'autant plus raide que $h$ est petit : en prenant $h = L_1$
        alors que la vraie zone de rupture est plus courte, on \emph{sous-estime} la sévérité du
        snap-back.
  \item La relation $G_f = e\,f_t^2L_1/E_1$ lie l'énergie de rupture à $L_1$. Si $L_1$ est la
        longueur de l'éprouvette et non celle de la zone de rupture, le $G_f$ induit est trop
        grand, ou, si l'on impose $G_f$, le $\uc$ qui en découle est trop petit pour la longueur
        d'éprouvette.
\end{itemize}

\subsection{Ce qu'elle épargne}
Avec une loi locale et $h$ libre, Bažant montre que le trajet stable est $h \to 0$ : dissipation
nulle, et résultat dépendant du maillage dans un calcul éléments finis (figure 13.43a : la
descente post-pic est de plus en plus raide quand le maillage se raffine, et l'énergie dissipée
tend vers zéro). Pour éviter cela il faut imposer que $h$ ne descende pas sous une longueur
caractéristique matériau : c'est le modèle de bande de fissuration (§13.10), où la largeur de
bande $w_c$ est une propriété du matériau, de l'ordre de quelques tailles de granulat, et où la
loi $\sigma(\varepsilon)$ est ajustée pour que $w_c\int\sigma\,\dd\varepsilon = G_f$.

Notre modèle est \emph{déjà} un modèle de bande de fissuration : $L_1$ joue le rôle de $w_c$, et
la relation $G_f = L_1\int\sigma\,\dd\varepsilon$ est exactement la sienne. C'est pour cela que le
modèle à un ressort est objectif sans aucune correction : la longueur de localisation y est
fixée par la physique (ou par un choix), pas par un maillage.

\begin{idee}
Le modèle à un ressort n'est pas une simplification arbitraire. C'est le modèle de bande de
fissuration de Bažant réduit à une seule bande, avec $w_c = L_1$. Il est objectif parce que
la longueur qui localise est imposée, et c'est précisément ce qu'il ne peut pas prédire.
\end{idee}

\section{Un ressort ou un segment ? La question de $L_1$}
La vraie question n'est donc pas « un ressort ou plusieurs », mais « que vaut $h$ ». Pour un
béton, la zone de rupture en traction a une largeur de quelques fois la taille du plus gros
granulat ($w_c \simeq 3d_a$ chez Bažant, soit quelques centimètres), et le reste de l'éprouvette
décharge. Si $L_1$ est pris égal à la longueur de l'éprouvette, on peut le corriger sans
changer la structure du modèle, en appliquant (13.2.5) :
\begin{itemize}[nosep]
  \item la zone qui s'adoucit a la longueur $h = w_c$ et la raideur tangente $C_a = E_tA_1/h$ ;
  \item le béton qui décharge, de longueur $L_1 - h$, s'ajoute en série à la barre d'acier, et
        la raideur de « tout le reste » devient
        \begin{equation}
          C_b = \Bigl(\frac1{k_2} + \frac{L_1 - h}{E_1A_1}\Bigr)^{-1} < k_2 .
        \end{equation}
\end{itemize}
Les deux effets vont dans le même sens : $C_a$ est plus raide en descente (division par $h$
au lieu de $L_1$), $C_b$ est plus souple. Le snap-back est plus sévère, et le point limite
plus proche du pic. Dans le code, cela revient à remplacer $L_1$ par $h$ dans la définition de
$k_1$ et de $\uc$, et $k_2$ par le $C_b$ ci-dessus, sans rien changer d'autre.

\begin{alamain}
Choisir une largeur de bande $h$ (par exemple $3d_a$ avec $d_a$ la taille du plus gros
granulat). Calculer $k_1' = E_1A_1/h$, $\uc' = \sqrt{2e}\,f_th/E_1$ (même $f_t$, $G_f' = ef_t^2h/E_1$),
et $C_b = (1/k_2 + (L_1-h)/(E_1A_1))^{-1}$. Vérifier le critère $C_b < 2e^{-3/2}k_1'$ et comparer
le point limite au cas $h = L_1$.
\end{alamain}

\section{Voir la localisation se produire : le modèle à $n$ ressorts}
Le moyen le plus parlant de relier les deux sens du mot est de découper le ressort 1 en $n$
ressorts identiques en série, chacun de longueur $L_1/n$, avec la même loi locale écrite en
déformation :
\begin{equation}
  \sigma_i = E_1(1-d_i)\,\varepsilon_i,\qquad \varepsilon_i = \frac{u_i - u_{i-1}}{L_1/n},\qquad
  \hat d_i = 1 - \exp\!\bigl(-\varepsilon_i^2/\varepsilon_c^2\bigr),\quad \varepsilon_c = \frac{\uc}{L_1}\ \text{fixé}.
\end{equation}
Le problème a alors $n$ inconnues $u_1,\ldots,u_n$ (le noeud $n$ est relié à la barre d'acier),
un résidu vectoriel $R_i = N_i - N_{i+1}$ (équilibre de chaque noeud, avec $N_{n+1} = k_2(u_2 - u_n)$)
et une matrice tangente tridiagonale. On observe alors, par le calcul, ce que Bažant prédit :
\begin{enumerate}[nosep]
  \item tant que $d = 0$, tous les segments sont identiques et la solution est uniforme ;
  \item au pic, le système est au point de bifurcation : l'état uniforme et les états
        localisés (un segment s'adoucit, $n-1$ déchargent le long de leur sécante) sont tous
        des équilibres. Dans un calcul, c'est l'arrondi ou une imperfection (un segment avec
        $f_t$ réduit de $0{,}1$~\%) qui choisit, et la matrice tangente de l'état uniforme a
        $n-1$ valeurs propres négatives, signe que ce trajet n'est pas le trajet stable ;
  \item la courbe $F(u_2)$ post-pic dépend de $n$ : plus raide à chaque raffinement, énergie
        dissipée en $1/n$, et le critère de snap-back se déplace,
        $k_2 < 2e^{-3/2}\,E_1A_1/(L_1/n)$, donc de plus en plus facile à satisfaire ;
  \item pour retrouver un résultat indépendant de $n$, il faut faire dépendre la loi de la
        taille du segment, $\varepsilon_c \propto \sqrt{n}$ de sorte que $G_f$ soit conservée : c'est
        la correction de bande de fissuration de Bažant (13.10.3), ou bien passer à un modèle
        non local (§13.10) ou à effet retard (Suffis et Combescure).
\end{enumerate}
C'est l'expérience numérique de la figure 13.43a de Bažant. Le modèle à un ressort apparaît
alors comme le cas $n = 1$ d'une famille, le seul qui soit objectif sans correction, parce
que $h$ y est fixé par la physique et non par la discrétisation.

\section{Lien avec la zone cohésive}
Une zone cohésive (Acary et Monerie, Chaboche et Feyel) pousse l'hypothèse de localisation à
sa limite : $h \to 0$ tout en conservant $G_f$ fini, en remplaçant la déformation par un saut
de déplacement $[u]$ et la contrainte par une traction $T([u])$. L'énergie est portée par une
surface, plus par un volume, donc le raffinement du maillage du volume ne change rien. Notre
ressort est à mi-chemin : une interphase d'épaisseur $L_1$ plutôt qu'une interface
d'épaisseur nulle (Chaboche et Feyel, §2.2). La différence pratique est que notre $u_1$
contient l'allongement élastique du béton intact en plus de l'ouverture de fissure, ce qui
rend la descente de $F_1(u_1)$ plus raide que celle d'une zone cohésive de mêmes $f_t$ et $G_f$ :
notre modèle est donc légèrement plus sujet au snap-back qu'une zone cohésive extrinsèque
équivalente (point 4 d'Acary et Monerie).

\begin{center}
\begin{tabular}{@{}p{3.4cm}p{3.6cm}p{3.6cm}p{3.6cm}@{}}
\toprule
 & Continuum local maillé & Notre ressort ($n = 1$) & Zone cohésive\\
\midrule
Zone qui localise & un élément, $h\to0$ & $h = L_1$, imposé & surface, $h = 0$\\
Énergie dissipée & $\to 0$ avec le maillage & $G_f = L_1\int\sigma\,\dd\varepsilon$ & $G_f = \int T\,\dd[u]$\\
Objectivité & non sans correction & oui, par construction & oui\\
Prédit où/combien localise & oui (mal) & non & non\\
Trajet stable au pic & bifurcation, descente la plus raide & pas de bifurcation & pas de bifurcation\\
Critère de snap-back & $C_b < -E_tA/h$, dépend de $h$ & $k_2 < 2e^{-3/2}k_1$ & $E/L < \sup|\xi'|$\\
\bottomrule
\end{tabular}
\end{center}

\part{La stabilité par le travail du second ordre}

\section{La phrase de Chaboche et Feyel}
Dans la section 3 de leur article (p.~3132), Chaboche, Feyel et Monerie écrivent du point
limite : « this bifurcation point is special: it induces a loss of stability, because the
second-order work is zero (Bažant and Cedolin, 1991), but the uniqueness is not lost ». Cette
phrase contient tout le sujet : le critère (le travail du second ordre), sa référence
(Bažant, §10.1 et §13.2), et une nuance essentielle (unicité conservée). On détaille les trois.

\section{Pourquoi le travail du second ordre}

\subsection{Le premier ordre ne dit rien}
Soit un système en équilibre sous charge fixée. Pour un incrément virtuel $\delta q$ des
déplacements, le travail des forces intérieures et extérieures au premier ordre est
\begin{equation}
  \delta\mathcal W = \sum_{\text{éléments}} F\,\delta(\text{allongement}) - P\,\delta q_{\text{charge}} .
\end{equation}
Le principe des travaux virtuels dit que $\delta\mathcal W = 0$ pour tout $\delta q$ admissible si
et seulement si le système est en équilibre. Chez nous, $\delta\mathcal W = R(u_1;u_2)\,\delta u_1$ :
le travail du premier ordre est le résidu fois l'incrément, et il est nul sur tout équilibre,
stable ou non. Il ne peut donc pas distinguer un équilibre stable d'un équilibre instable.

\subsection{Le second ordre}
Le premier terme non nul du développement de l'énergie autour d'un équilibre est le terme
d'ordre deux. En notant $\delta f$ l'incrément de force qui accompagne $\delta q$ le long de la
réponse tangente, le travail du second ordre est
\begin{equation}
  \dW = \tfrac12\,\delta f\,\delta q = \tfrac12\,\delta q^{\mathsf T}\,K^t\,\delta q ,
\end{equation}
avec $K^t$ la matrice de raideur tangente. C'est l'énergie qu'il faudrait \emph{fournir} de
l'extérieur pour imposer l'incrément $\delta q$ à partir de l'équilibre. Bažant (§10.1) lui donne
un sens thermodynamique :
\begin{equation}
  \dW = -T\,(\Delta S)_{\text{in}} ,
\end{equation}
l'opposé de l'entropie produite par la transformation (à température $T$, en conditions
isothermes ; c'est la variation seconde de l'énergie libre de Helmholtz du système
structure + chargement). Le critère en découle par le second principe :
\begin{itemize}[nosep]
  \item si $\dW > 0$ pour tout $\delta q$ admissible, toute excursion coûte de l'énergie et
        produit une entropie négative, ce qui est impossible spontanément : l'état est
        \textbf{stable} ;
  \item si $\dW < 0$ pour un $\delta q$ admissible, cette excursion libère de l'énergie et produit
        de l'entropie ; le second principe impose qu'elle se produise : l'état est
        \textbf{instable}, et l'énergie libérée part en énergie cinétique puis, par
        l'amortissement inévitable, en chaleur ;
  \item si $\dW = 0$ pour un $\delta q$ et $> 0$ pour les autres, l'état est \textbf{critique}
        (équilibre neutre au premier ordre non nul) ; la stabilité est alors décidée par les
        variations d'ordre supérieur.
\end{itemize}
C'est le théorème de Lagrange--Dirichlet étendu aux systèmes non élastiques : pour un système
élastique, $\dW$ est la variation seconde de l'énergie potentielle et le critère est celui du
minimum d'énergie ; pour un système avec dissipation, Bažant montre que l'entropie produite
en est l'extension correcte.

\subsection{Le mot « admissible »}
Le critère porte sur les incréments \emph{admissibles}, c'est-à-dire compatibles avec les
conditions aux limites et avec les lois de comportement en charge et en décharge. Deux
conséquences, toutes deux visibles sur notre système.
\begin{itemize}[nosep]
  \item Les conditions aux limites restreignent $\delta q$. En \emph{force} pilotée, le
        déplacement du point de chargement est libre ; en \emph{déplacement} piloté, il est
        bloqué ($\delta q_1 = 0$). Le même état peut donc être instable en force et stable en
        déplacement : c'est tout l'adoucissement entre le pic et le point limite.
  \item Les lois de comportement font que $K^t$ dépend du sens de l'incrément : en charge la
        raideur tangente de notre ressort 1 est $\kd + u_1\,\dd\kd/\dd u_1$, en décharge c'est
        $\kd$ seulement. Il faut vérifier $\dW > 0$ pour toutes les combinaisons de charge et
        de décharge des différentes parties : Bažant insiste, « one needs to check various
        combinations of loading and unloading stiffnesses for stability ».
\end{itemize}

\section{Application à notre système}

\subsection{Les deux parties}
Dans le langage de Bažant (§13.2), partie $a$ = ressort 1, raideur tangente $C_a$ ; partie
$b$ = ressort 2, raideur tangente $C_b = k_2$ ; $q_2 = u_1$ le déplacement de l'interface,
$q_1 = u_2$ le déplacement du point de chargement. Les incréments de force sont
$\delta F_a = C_a\,\delta u_1$ et $\delta F_b = k_2(\delta u_2 - \delta u_1)$, et
\begin{equation}
  \dW = \tfrac12\,C_a\,\delta u_1^2 + \tfrac12\,k_2\,(\delta u_2 - \delta u_1)^2 .
  \tag{13.2.1}
\end{equation}
Pour $C_a$, deux valeurs selon le sens de l'incrément :
\begin{equation}
  C_a^{\text{charge}} = \frac{\dd F_1}{\dd u_1} = k_1e^{-x^2}(1-2x^2),\qquad
  C_a^{\text{décharge}} = \kd = k_1e^{-x^2} > 0 .
\end{equation}

\subsection{Pilotage en force}
$\delta u_1$ et $\delta u_2$ sont tous deux libres. La forme quadratique (13.2.1) est définie
positive si et seulement si $C_a > 0$ et $k_2 > 0$. En décharge c'est toujours vrai ; en charge
il faut $\dd F_1/\dd u_1 > 0$, donc $u_1 < \uc/\sqrt2$. Au-delà du pic, l'incrément « $\delta u_1 > 0$
en charge, $\delta u_2 = \delta u_1 + C_a\delta u_1/k_2$ qui maintient la force » donne $\dW < 0$ :
instable. On ne peut pas dépasser le pic en force, quel que soit $k_2$.

\subsection{Pilotage en déplacement}
$\delta u_2 = 0$ pour toute excursion possible, et (13.2.1) devient
\begin{equation}
  \dW = \tfrac12\,(C_a + k_2)\,\delta u_1^2 .
  \tag{13.2.2}
\end{equation}
En décharge, $C_a + k_2 > 0$ toujours. En charge, $C_a + k_2 = k_t$, et le critère est
\begin{equation}
  \text{stable si } k_t > 0,\qquad \text{critique si } k_t = 0,\qquad \text{instable si } k_t < 0 .
  \tag{13.2.3}
\end{equation}
C'est la condition $-C_a < C_b$ de Bažant. Elle est atteinte au point limite, qui est donc l'état
critique au sens du travail du second ordre, exactement ce que dit Chaboche. Bažant en donne une
lecture directe : la force qu'il faudrait appliquer à l'interface pour maintenir l'équilibre
après un déplacement $\delta u_1$ est $\delta f_2 = (C_a + k_2)\,\delta u_1 = k_t\,\delta u_1$ ; si son
travail $\tfrac12\delta f_2\,\delta u_1$ est positif, aucun changement d'état ne peut se produire
sans qu'on applique cette force ; s'il est négatif, l'état se défait tout seul.

\subsection{L'asymétrie charge--décharge}
Sur un état de la branche descendante où $k_t < 0$ en charge, le même état donne
$C_a^{\text{décharge}} + k_2 = \kd + k_2 > 0$ en décharge. L'état est donc instable pour
$\delta u_1 > 0$ et stable pour $\delta u_1 < 0$ : il se défait en s'ouvrant, jamais en se
refermant. C'est ce qui garantit qu'après le saut la masse oscille sans que l'endommagement
ne reparte, et c'est pourquoi la surface $\dW(\delta q_1,\delta q_2)$ de Bažant (figure 13.7)
est faite de secteurs quadratiques raccordés avec continuité de pente mais discontinuité de
courbure.

\subsection{L'interprétation énergétique de Bažant}
Pour une chute de force $\delta P$ le long de la réponse, la décharge de $b$ libère
$(\delta P)^2/2k_2$ et l'adoucissement de $a$ requiert $-(\delta P)^2/2C_a$. Si
$(\delta P)^2/2k_2 < -(\delta P)^2/2C_a$, soit $k_2 > -C_a$, l'énergie libérée ne suffit pas à
nourrir l'adoucissement : stable. Sinon « there is an excess release of energy that can
manifest itself as nothing else but kinetic energy, and so equilibrium can no longer exist ».
Intégré sur tout le saut, cet excès est notre $|W|$.

\section{La nuance de Chaboche : stabilité perdue, unicité conservée}
\label{sec:trajet}

\subsection{Bifurcation et point limite}
Bažant distingue deux situations où $\dW$ s'annule.
\begin{itemize}[nosep]
  \item Au \emph{point de bifurcation}, plusieurs trajets d'équilibre partent du même état
        (flambement d'Euler, colonne de Shanley, localisation au pic dans une barre maillée).
        L'état lui-même peut être stable ($\dW > 0$ pour toute excursion à charge fixée) ; ce qui
        est en jeu est le choix du trajet, et le critère est différent : le trajet suivi est
        celui qui minimise $\dW$ le long du trajet, c'est-à-dire, à déplacement imposé, celui de
        plus petite raideur globale $C$, « the steepest downward slope » (§10.2 et §13.2).
  \item Au \emph{point limite}, un seul trajet existe, mais il rebrousse chemin par rapport au
        paramètre de contrôle ; $\dW = 0$ parce que la raideur de l'état devient nulle dans la
        direction du trajet. Il n'y a pas de choix : l'état perd sa stabilité et le système
        quitte le trajet d'équilibre.
\end{itemize}
Notre point limite est du second type. C'est pourquoi Chaboche écrit que l'unicité n'est pas
perdue : la branche d'équilibre reste unique (c'est notre courbe analytique), et l'état rompu
est l'unique solution au-delà, « instantaneous failure is obviously the exact solution of the
present elementary mechanical problem ». Ce qui est perdu est la stabilité de l'état, et avec
elle la continuité de la solution par rapport à $u_2$ croissant. Bažant insiste sur la même
distinction au §13.2 : la bifurcation avec localisation survient au pic si la localisation est
possible, l'instabilité par snap-back survient plus tard, « at a certain finite value of
negative tangent modulus ». Dans notre modèle à un degré de liberté la première n'existe pas,
seule la seconde reste.

\subsection{Pourquoi les itérations de Newton ne voient que la stabilité de l'état}
Bažant (p.~909) : « the convergence of these iterations has nothing to do with path stability.
It depends only on the stability of state. The reason is that, during the iterations of a
loading step, the loading does not change. Consequently, the errors in the iterated states
represent deviations away from the equilibrium state. » Autrement dit, Newton teste $\dW$ à
chargement fixé, donc le critère d'état (13.2.3), et rien d'autre. Sur notre système, Newton
converge tant que $k_t > 0$ et échoue au point limite : il mesure exactement la perte de
stabilité de l'état. Dans un système à plusieurs degrés de liberté, il convergerait tout aussi
bien sur un trajet instable (déformation uniforme au-delà du pic) tant que chaque état y est
stable à chargement fixé ; le seul moyen de détecter le mauvais trajet est alors de regarder
le signe des valeurs propres (ou des pivots) de $K^t$, ce qui est en 1D notre vérification
$k_t > 0$ à convergence.

\section{Résumé des critères}
\begin{center}
\begin{tabular}{@{}p{3.3cm}p{5.0cm}p{5.6cm}@{}}
\toprule
Question & Critère & Chez nous\\
\midrule
Équilibre & $\delta\mathcal W = 0$ pour tout $\delta q$ & $R(u_1;u_2) = 0$\\
Stabilité de l'état & $\dW = \tfrac12\delta q^{\mathsf T}K^t\delta q > 0$ pour tout $\delta q$ admissible (charge et décharge) & force pilotée : $C_a > 0$ (avant le pic) ; déplacement piloté : $k_t > 0$ (avant le point limite)\\
État critique & $\dW = 0$ pour un $\delta q$ & $k_t = 0$, point limite, tangente verticale de $F(u_2)$\\
Trajet stable (bifurcation) & minimiser $\dW$ le long du trajet, plus petit $C$ & sans objet à un degré de liberté ; localisation au pic si l'on maille\\
Unicité & & conservée au point limite (Chaboche) : un seul trajet, qui rebrousse\\
Ce que mesure Newton & stabilité de l'état à chargement fixé & converge si $k_t > 0$, échoue sinon\\
\bottomrule
\end{tabular}
\end{center}

\begin{alamain}
Calculer $C_a^{\text{charge}}(u_1)$ et $C_a^{\text{décharge}}(u_1)$ le long de la branche
analytique ; tracer $\dW/\delta u_1^2 = \tfrac12(C_a + k_2)$ dans les deux sens ; repérer où il
s'annule en charge (point limite) et vérifier qu'il reste positif en décharge partout.
Calculer la raideur globale $C = (C_a^{-1} + k_2^{-1})^{-1}$ et vérifier qu'elle passe par
$-\infty$ puis $+\infty$ au point limite.
\end{alamain}

\begin{idee}
Le travail du premier ordre dit si l'on est à l'équilibre ; le travail du second ordre dit si
cet équilibre tient. Pour notre système en déplacement piloté, $\dW = \tfrac12k_t\,\delta u_1^2$ :
le signe de la raideur tangente décide, et le point limite est l'état critique. L'unicité
n'est pas en cause, la stabilité l'est.
\end{idee}

\section*{Références}
\begin{itemize}[nosep]
  \item Z.~P. Bažant, L. Cedolin, \emph{Stability of Structures}, 1991 : §10.1 (travail du second
        ordre et entropie), §10.2 (trajet stable), §13.2 (couplage en série, équations 13.2.1 à
        13.2.5, figure 13.7), §13.8 p.~909 (convergence des itérations), §13.9--13.10 (sensibilité
        au maillage, bande de fissuration).
  \item J.-L. Chaboche, F. Feyel, Y. Monerie, IJSS 2001, §3 : « loss of stability because the
        second-order work is zero, but the uniqueness is not lost » ; §2.2 interphase/interface.
  \item R. de Borst, C. Verhoosel, 2017 : critère de Hill, barre à $m$ éléments.
  \item V. Acary, Y. Monerie, 2006 : zone cohésive, théorème d'unicité, relaxation dynamique.
  \item Note de cours \emph{Loi de comportement élastique-endommageable} : poutre non uniforme.
\end{itemize}

\end{document}
